\section{第02章 ソフトウェアアーキテクチャ}
\noindent この資料の主張\par
ソフトウェアアーキテクチャは、単なる「大まかな設計図」ではない。アーキテクチャは、ソフトウェア集約システムにおいて根本的であり、後から変えにくく、多くの品質や費用に影響する構造・関係・原則・重要な意思決定の集合である。よいアーキテクチャは、利害関係者の関心事を整理し、設計・実装・テスト・運用・保守に関わる人々が同じ理解を持つための土台になる。

\begin{statementbox}{この資料の読み方}
専門用語は原則として日本語で示し、必要に応じて英語を括弧内に併記する。英語名は、元資料、規格、ツール、文献で同じ概念を確認するための手がかりとして残す。初学者が前提知識なしに読めるように、各節では「何のための概念か」「実務では何を判断するのか」を重視する。
\end{statementbox}
